\begin{appendices}
\chapter{List of Symbols and Abbrevations}


% Add these packages to your document preamble:
% \usepackage{longtable}
% \usepackage{booktabs}
% \usepackage{array}

%========================================================
% APPENDIX A: LIST OF SYMBOLS AND ABBREVIATIONS
%========================================================


%--------------------------------------------------------
\section*{Latin Symbols}
%--------------------------------------------------------

\renewcommand{\arraystretch}{1.15}
\setlength{\tabcolsep}{5pt}

\begin{longtable}{@{}p{0.17\textwidth}
                      p{0.63\textwidth}
                      p{0.15\textwidth}@{}}
\toprule
\textbf{Symbol} & \textbf{Description} & \textbf{Unit} \\
\midrule
\endfirsthead

\toprule
\textbf{Symbol} & \textbf{Description} & \textbf{Unit} \\
\midrule
\endhead

\bottomrule
\endfoot

$A$ & Sectional or waterplane area, as applicable & m$^2$ \\
$A_{\mathrm{midships}}$ & Immersed midship sectional area & m$^2$ \\
$A_{\mathrm{tunnel}}$ & Tunnel surface area allowance & m$^2$ \\
$A_{WP}$ & Waterplane area & m$^2$ \\

$B$ & Moulded breadth of the vessel & m \\
$B_{WL}$ & Breadth at the waterline & m \\
$BM$ & Metacentric radius & m \\
$BM_L$ & Longitudinal metacentric radius & m \\
$BM_T$ & Transverse metacentric radius & m \\

$C_A$ & Correlation allowance coefficient & -- \\
$C_{AA}$ & Air resistance coefficient, as defined in the resistance calculation & -- \\
$C_{AS}$ & Steering resistance coefficient & -- \\
$C_B$ & Block coefficient & -- \\
$C_F$ & Frictional resistance coefficient & -- \\
$C_M$ & Midship section coefficient & -- \\
$C_P$ & Prismatic coefficient & -- \\
$C_R$ & Residual resistance coefficient & -- \\
$C_T$ & Total resistance coefficient & -- \\
$C_W$ & Waterplane area coefficient & -- \\

$d$ & Centre-girder depth in the structural rule calculation & mm \\
$D$ & Moulded depth of the vessel & m \\

$Fn$ & Froude number & -- \\
$G$ & Centre of gravity & Reference point \\
$g$ & Acceleration due to gravity & m/s$^2$ \\
$GB$ & Distance between the centres of gravity and buoyancy & m \\
$GM$ & Metacentric height & m \\
$GM_L$ & Longitudinal metacentric height & m \\
$GM_T$ & Transverse metacentric height & m \\
$GZ$ & Righting arm & m \\

$h$ & Spacing between successive ordinates or waterlines & m \\
$I_L$ & Longitudinal second moment of waterplane area & m$^4$ \\
$I_T$ & Transverse second moment of waterplane area & m$^4$ \\

$K$ & Keel reference point & Reference point \\
$KB$ & Vertical distance from keel to centre of buoyancy & m \\
$KG$ & Vertical distance from keel to centre of gravity & m \\
$KM$ & Height of metacentre above keel & m \\
$KM_L$ & Longitudinal metacentric height above keel & m \\
$KM_T$ & Transverse metacentric height above keel & m \\

$L$ & Reference length of the vessel & m \\
$L/B$ & Length-to-breadth ratio & -- \\
$L/T$ & Length-to-draft ratio & -- \\
$L_{BP}$ & Length between perpendiculars & m \\
$L_{OA}$ & Length overall & m \\
$L_{WL}$ & Length at the waterline & m \\

$LCB$ & Longitudinal centre of buoyancy & m \\
$LCF$ & Longitudinal centre of flotation & m \\

$M$ & Moment ordinate used in numerical integration & Context-dependent \\
$M_{BW}$ & Mass of ballast water & t \\
$M_{DF}$ & Mass of diesel fuel & t \\
$M_{FW}$ & Mass of fresh water & t \\
$M_{LO}$ & Mass of lubricating oil & t \\
$M_{PR}$ & Mass of provisions & t \\
$MCT$ & Moment to change trim by one centimetre & t\,m/cm \\

$N_{\mathrm{crew}}$ & Number of crew members & persons \\

$P_B$ & Brake power & W \\
$P_D$ & Delivered power & W \\
$P_E$ & Effective power & W \\

$QPC$ & Quasi-propulsive coefficient & -- \\

$R_T$ & Total resistance & N \\
$Re$ & Reynolds number & -- \\

$S$ & Wetted surface area of the hull & m$^2$ \\

$T$ & Draft of the vessel & m \\
$TPC$ & Tonnes per centimetre immersion & t/cm \\

$V$ & Ship speed in resistance and powering equations & m/s \\
$V_{CB}$ & Vertical centre of buoyancy coordinate & m \\
$V_{md}$ & Volume up to the main deck & m$^3$ \\

$W$ & Weight or mass allowance, according to context & Context-dependent \\

$y_0$ & Half-ordinate at the reference waterline & m \\


\end{longtable}

%--------------------------------------------------------
\section*{Greek Symbols}
%--------------------------------------------------------

\begin{longtable}{@{}p{0.17\textwidth}
                      p{0.63\textwidth}
                      p{0.15\textwidth}@{}}
\toprule
\textbf{Symbol} & \textbf{Description} & \textbf{Unit} \\
\midrule
\endfirsthead

\toprule
\textbf{Symbol} & \textbf{Description} & \textbf{Unit} \\
\midrule
\endhead

\bottomrule
\endfoot

$\Delta$ & Mass displacement of the vessel & t \\
$\Delta C_R$ & Correction to the residual resistance coefficient & -- \\
$\Delta T$ & Change in draft & m \\
$\Delta W$ & Change in weight & t \\
$\Delta V$ & Change in submerged volume & m$^3$ \\

$\eta$ & Propulsion efficiency & -- \\

$\theta$ & Angle of heel & degrees or rad \\
$\nabla$ & Volume of displacement & m$^3$ \\

$\nu$ & Kinematic viscosity of water & m$^2$/s \\
$\rho$ & Density of water & kg/m$^3$ \\

\end{longtable}

%--------------------------------------------------------
\section*{Abbreviations}
%--------------------------------------------------------

\begin{longtable}{@{}p{0.22\textwidth}
                      p{0.72\textwidth}@{}}
\toprule
\textbf{Abbreviation} & \textbf{Full Form} \\
\midrule
\endfirsthead

\toprule
\textbf{Abbreviation} & \textbf{Full Form} \\
\midrule
\endhead

\bottomrule
\endfoot


AP & After Perpendicular \\
ASW & Anti-Submarine Warfare \\
B/T & Breadth-to-draft ratio \\
CAD & Computer-Aided Design \\
CFD & Computational Fluid Dynamics \\
CODAG & Combined Diesel and Gas Turbine \\
COG & Centre of Gravity \\
DE & Diesel-Electric \\
FP & Forward Perpendicular \\
IMO & International Maritime Organization \\
IRS & Indian Register of Shipping \\
ISO & International Organization for Standardization \\
LCB & Longitudinal Centre of Buoyancy \\
LCF & Longitudinal Centre of Flotation \\
LBP & Length Between Perpendiculars \\
LOA & Length Overall \\
LWL & Length at Waterline \\
MCT & Moment to Change Trim by One Centimetre \\
NPL & National Physical Laboratory (hull series) \\
QPC & Quasi-Propulsive Coefficient \\
TPC & Tonnes per Centimetre Immersion \\
VLS & Vertical Launching System \\

\end{longtable}

%========================================================
% APPENDIX B: GLOSSARY OF TERMS
%========================================================
\clearpage
\thispagestyle{empty}
\null
\newpage
\chapter{Glossary of Terms}

\begin{description}[
    leftmargin=4.3cm,
    labelwidth=3.8cm,
    labelsep=0.25cm,
    align=left,
    style=multiline,
    itemsep=0.65em,
    parsep=0pt,
    font=\normalfont\bfseries
]

\item[Appendages]
Components projecting from or attached to the main hull, such as rudders, bilge keels, shaft brackets, and stabilising fins, which may affect resistance and manoeuvring.

\item[Bending Moment]
The internal moment produced by the distribution of loads along a structural member. In a ship's hull girder, it causes longitudinal bending under the combined effects of weight and buoyancy.

\item[Cavitation]
The formation of vapour cavities in a liquid when local pressure falls sufficiently low, followed by their collapse when pressure increases. In propellers, it can cause noise, vibration, erosion, and loss of performance.

\item[Complement]
The total number of personnel required to operate, maintain, and support a ship.

\item[Course-Keeping]
The ability of a ship to maintain a desired course with minimal deviation despite environmental disturbances.

\item[Freeboard]
The vertical distance measured at the ship's side between the waterline and the applicable freeboard deck line or deck reference.

\item[Froude Number ($Fn$)]
A dimensionless number expressing the ratio of inertial effects to gravitational effects in free-surface flow. It is given by
$Fn = \frac{V}{\sqrt{gL}}$, where $V$ is ship speed, $g$ is gravitational acceleration, and $L$ is a characteristic length.

\item[General Arrangement (GA)]
A drawing showing the layout of a ship's compartments, decks, machinery, equipment, accommodation, and other principal spaces.

\item[Length of Run ($L_R$)]
The longitudinal length of the afterbody over which the hull narrows towards the stern and the flow is directed towards the propeller. Its precise limits depend on the hull-form convention used.

\item[Lines Plan]
A set of technical drawings representing the three-dimensional hull form through the body plan, half-breadth plan, and sheer plan.

\item[Maneuverability]
The ability of a ship to change its speed, heading, or direction of motion in response to control inputs.

\item[Metacentric Height ($GM$)]
The vertical distance between the ship's centre of gravity and metacentre. It is an important measure of the ship's initial static stability.

\item[Metacentric Radius ($BM$)]
The distance between the centre of buoyancy and the metacentre. It is calculated as the waterplane's second moment of area about the relevant axis divided by the underwater displacement volume.

\item[Offset Table]
A table containing the coordinates of the hull surface at specified stations and waterlines, used to generate and represent the ship's hull geometry.

\item[Payload]
The total weight or mass of mission-specific equipment, weapons, sensors, cargo, and other operational loads carried by a vessel.

\item[Quasi-Propulsive Coefficient (QPC)]
The ratio of effective power required to overcome ship resistance to the delivered power supplied to the propeller, accounting for hull--propeller interaction effects.

\item[Redundancy]
The provision of duplicate or alternative equipment and systems so that essential functions can continue if a component fails.

\item[Reynolds Number ($Re$)]
A dimensionless number representing the ratio of inertial forces to viscous forces in a fluid flow. It is given by
$Re = \frac{VL}{\nu}$, where $V$ is characteristic velocity, $L$ is characteristic length, and $\nu$ is kinematic viscosity.

\item[Reserve Buoyancy]
The watertight volume of a ship above the current waterline that can provide additional buoyancy as the vessel sinks deeper into the water, provided the relevant boundaries remain watertight.

\item[Seakeeping]
The ability of a ship to maintain safe and effective operation in waves, considering its motions, accelerations, and overall performance.

\item[Shear Force]
The internal force acting across a structural section due to the distribution of external loads. In a ship's hull girder, it results from the longitudinal distribution of weight and buoyancy.

\item[Sheer]
The longitudinal curvature or rise of a ship's deck towards the bow and stern relative to the deck line at midship.

\item[Stealth Design]
The application of design techniques and materials to reduce a ship's detectability by radar, infrared, acoustic, and other sensing systems.

\item[Superstructure]
The structure built above the main deck of a ship, generally accommodating navigation spaces, operational equipment, and other facilities.

\item[Turbine]
A rotary machine that extracts energy from a moving fluid and converts it into mechanical power. In marine propulsion, gas turbines may be used to drive the propulsion system.

\end{description}

% ============================================================
% APPENDIX C - SHIP MAIN DIMENSION GENERATOR
% ============================================================

\chapter{Python code used for  Estimation of Dimensions}
\label{app:ship_dimension_generator}

The following Python program was developed to generate feasible
combinations of the principal dimensions of the proposed vessel.
The program evaluates combinations of length overall (LOA), breadth
(B), depth (D), block coefficient ($C_B$), and displacement, and
calculates the corresponding waterline length, breadth at the design
waterline, draft, and principal dimensional ratios.

\begin{lstlisting}[language=Python,
    caption={Python program for generating feasible ship main dimensions},
    label={lst:ship_dimension_generator}]
import csv

# ============================================================
# SHIP MAIN DIMENSION GENERATOR
# ============================================================

# MAIN DIMENSIONS
LOA_MIN = 114
LOA_MAX = 118
LOA_STEP = 1

B_MIN = 14
B_MAX = 16
B_STEP = 1

D_MIN = 9
D_MAX = 10
D_STEP = 1

# BLOCK COEFFICIENT
CB_VALUES = [0.380, 0.39, 0.395, 0.400]

# DISPLACEMENT
DISP_MIN = 3600
DISP_MAX = 3850
DISP_STEP = 50

# CONSTANTS
LWL_RATIO = 0.975
B_DWL_RATIO = 0.94
RHO = 1.025


# ============================================================
# GENERATE INTEGER VALUES
# ============================================================

LOA_VALUES = range(LOA_MIN, LOA_MAX + 1, LOA_STEP)
B_VALUES = range(B_MIN, B_MAX + 1, B_STEP)
D_VALUES = range(D_MIN, D_MAX + 1, D_STEP)
DISP_VALUES = range(DISP_MIN, DISP_MAX + 1, DISP_STEP)


# ============================================================
# CALCULATE POSSIBLE DIMENSIONS
# ============================================================

results = []

for LOA in LOA_VALUES:

    # LWL is calculated from LOA
    LWL = 0.975 * LOA

    for B in B_VALUES:

        # Breadth at DWL
        B_DWL = 0.94 * B

        for D in D_VALUES:

            for CB in CB_VALUES:

                for displacement in DISP_VALUES:

                    # ----------------------------------------
                    # DISPLACEMENT EQUATION
                    #
                    # Delta = rho x LWL x B_DWL x T x CB
                    #
                    # T = Delta / (rho x LWL x B_DWL x CB)
                    # ----------------------------------------

                    T = displacement / (
                        RHO *
                        LWL *
                        B_DWL *
                        CB
                    )

                    # Ratios
                    L_B = LOA / B
                    B_T = B / T
                    D_T = D / T

                    # Draft must be less than depth
                    if T < D:

                        results.append([
                            LOA,
                            LWL,
                            B,
                            B_DWL,
                            D,
                            CB,
                            displacement,
                            T,
                            L_B,
                            B_T,
                            D_T
                        ])


# ============================================================
# DISPLAY RESULTS
# ============================================================

print()
print("=" * 120)
print("POSSIBLE SHIP MAIN DIMENSIONS")
print("=" * 120)

print()

print(
    f"{'No.':<5}"
    f"{'LOA':<7}"
    f"{'LWL':<9}"
    f"{'B':<7}"
    f"{'B_DWL':<9}"
    f"{'D':<7}"
    f"{'CB':<8}"
    f"{'Disp':<9}"
    f"{'Draft':<9}"
    f"{'L/B':<8}"
    f"{'B/T':<8}"
    f"{'D/T':<8}"
)

print("-" * 120)


for i, row in enumerate(results, 1):

    (
        LOA,
        LWL,
        B,
        B_DWL,
        D,
        CB,
        displacement,
        T,
        L_B,
        B_T,
        D_T
    ) = row

    print(
        f"{i:<5}"
        f"{LOA:<7}"
        f"{LWL:<9.3f}"
        f"{B:<7}"
        f"{B_DWL:<9.2f}"
        f"{D:<7}"
        f"{CB:<8.3f}"
        f"{displacement:<9}"
        f"{T:<9.3f}"
        f"{L_B:<8.2f}"
        f"{B_T:<8.2f}"
        f"{D_T:<8.2f}"
    )


# ============================================================
# SAVE RESULTS TO CSV
# ============================================================

filename = "ship_possible_dimensions.csv"

with open(
    filename,
    "w",
    newline="",
    encoding="utf-8"
) as file:

    writer = csv.writer(file)

    writer.writerow([
        "LOA (m)",
        "LWL (m)",
        "Breadth B (m)",
        "Breadth at DWL (m)",
        "Depth D (m)",
        "CB",
        "Displacement (t)",
        "Draft T (m)",
        "L/B",
        "B/T",
        "D/T"
    ])

    for row in results:
        writer.writerow(row)


# ============================================================
# FINAL MESSAGE
# ============================================================

print()
print("=" * 120)
print("TOTAL FEASIBLE COMBINATIONS =", len(results))
print("=" * 120)
print()
print("Results saved as:", filename)
\end{lstlisting}


\end{appendices}